\section{Pengantar Permasalahan}

Kebijakan moneter merupakan salah satu instrumen utama negara dalam menjaga stabilitas ekonomi. Di Indonesia, kewajiban Bank Indonesia (BI) sebagai otoritas moneter berakar dari UU No. 23 Tahun 1999 tentang Bank Indonesia yang telah diubah beberapa kali, dengan perubahan terakhir melalui UU No. 4 Tahun 2023 tentang Pengembangan dan Penguatan Sektor Keuangan (UU P2SK). Pasal 9 angka 2 UU P2SK merumuskan tujuan Bank Indonesia sebagai berikut.

\begin{quote}
    \textit{Tujuan Bank Indonesia adalah mencapai stabilitas nilai rupiah, memelihara stabilitas Sistem Pembayaran, dan turut menjaga Stabilitas Sistem Keuangan dalam rangka mendukung pertumbuhan ekonomi yang berkelanjutan.}
\end{quote}

Stabilitas nilai rupiah dalam mandat tersebut secara eksplisit mencakup kestabilan harga barang dan jasa (inflasi) serta kestabilan nilai tukar rupiah terhadap mata uang negara lain, yang operasionalnya diatur lebih lanjut dalam Peraturan Bank Indonesia (PBI) No. 11 Tahun 2024 tentang Pengendalian Moneter yang menetapkan suku bunga acuan BI-Rate sebagai sasaran operasional utama kebijakan moneter. Pada 4 Juni 2026, DPR RI mengesahkan revisi lanjutan UU P2SK yang memperluas mandat BI untuk turut mendukung pertumbuhan ekonomi dan penciptaan lapangan kerja di samping tugas stabilisasi yang sudah ada, sehingga kebijakan moneter BI tidak dapat lagi dirancang semata-mata untuk menekan inflasi tetapi harus mempertimbangkan dampaknya terhadap sektor riil dan daya beli masyarakat.

Mandat hukum di atas secara langsung menghasilkan sebuah \textit{trade-off} multi-objektif yang harus diselesaikan Gubernur BI setiap periode Rapat Dewan Gubernur (RDG). Menaikkan suku bunga akan menekan inflasi dan memperkuat rupiah, tetapi menekan \textit{output} dan daya beli serta memberatkan biaya utang pemerintah. Sebaliknya, menurunkan suku bunga akan mendorong \textit{output} dan daya beli, tetapi berisiko melemahkan rupiah dan memicu \textit{capital outflow}.

\begin{figure}[H]
    \centering
    \includegraphics[width=0.5\textwidth]{rdg-bi.png}
    \caption{Rapat Dewan Gubernur Bank Indonesia. Sumber}
    \label{fig:rdg-bi}
\end{figure}

Sebagai ilustrasi kondisi riil yang dihadapi, pada RDG Juli 2026 BI mempertahankan BI-Rate di 5,75\% di tengah inflasi tahunan 3,34\% (naik dari 3,08\% bulan sebelumnya), nilai tukar rupiah di sekitar Rp17.885 per USD yang melemah 7,29\% \textit{year-to-date}, serta Indonesia baru mengalami defisit transaksi berjalan pertama setelah 72 bulan surplus berturut-turut. Kondisi ini terjadi ketika Federal Reserve AS menahan Fed Funds Rate di kisaran 3,50\% sampai 3,75\%, sehingga selisih suku bunga BI-Rate terhadap suku bunga AS menjadi pertimbangan penting untuk mencegah arus keluar modal asing.

Dari uraian di atas, dapat terlihat bahwa keputusan suku bunga tidak berdiri sendiri per periode dan dampaknya menjalar melalui inflasi, \textit{output}, dan nilai tukar ke periode-periode berikutnya. Menentukan satu angka BI-Rate untuk satu kuartal saja sudah cukup sulit karena melibatkan \textit{trade-off} yang dijelaskan sebelumnya. Menentukan delapan angka sekaligus untuk dua tahun ke depan, dengan setiap keputusan saling memengaruhi melalui rantai kausal yang kompleks, jelas berada di luar kemampuan intuisi atau \textit{rule of thumb} sederhana. Pendekatan naif seperti mematok suku bunga konstan atau mengikuti pola historis akan mengabaikan dinamika antar-periode yang non-linear serta interaksi antara enam \textit{constraint} yang harus dipenuhi secara bersamaan.

Maka dari itu, permasalahan ini akan dicoba untuk direduksi menjadi sebuah persoalan \textit{local search} yang mencari \textit{complete-state} berupa konfigurasi suku bunga sepanjang delapan kuartal yang memaksimalkan kesejahteraan sosial. Dengan formulasi ini, penentuan BI-Rate tidak lagi bergantung pada intuisi atau penilaian subjektif pengambil kebijakan, melainkan dijalankan secara sistematis melalui eksplorasi ruang state yang dipandu oleh fungsi objektif kuantitatif. Harapannya, pendekatan ini dapat menjadi kerangka bantu keputusan yang lebih terstruktur dibandingkan metode konvensional yang bersifat \textit{ad-hoc}.

\section{Tujuan Permasalahan}

Tujuan utama permasalahan ini adalah menentukan state suku bunga acuan $s = (r_1, r_2, \dots, r_T)$ yang memaksimalkan kesejahteraan gabungan dari stabilitas harga, stabilitas nilai tukar, dan daya beli masyarakat. Horizon kebijakan yang digunakan adalah $T = 8$ periode kuartalan (dua tahun ke depan), dimulai dari kondisi aktual BI-Rate per Juli 2026. Tiap $r_t$ harus berada dalam rentang 3,00\% sampai 8,00\% dengan kelipatan 25 basis poin sesuai konvensi RDG BI. State $s$ harus memenuhi beberapa \textit{constraint} yang mencerminkan kondisi neraca keuangan riil Indonesia. Secara formal, permasalahan dinyatakan sebagai pencarian state optimal pada ruang keadaan $\mathcal{S}$.

\[
    s^{*} = \arg\max_{s \,\in\, \mathcal{S}} J(s)
\]

Fungsi objektif $J(s)$ merupakan negatif dari \textit{quadratic loss function} yang mengukur total deviasi variabel makroekonomi dari targetnya sepanjang horizon $T$. Semakin mendekati nol nilai $J(s)$, semakin optimal state tersebut dalam menyeimbangkan inflasi, \textit{output}, daya beli, dan stabilitas suku bunga.

\section{Kerangka Teoretis Model Makroekonomi}

Untuk mengevaluasi seberapa baik sebuah state suku bunga $s = (r_1, \dots, r_T)$, diperlukan model yang mampu mensimulasikan bagaimana keputusan BI-Rate pada tiap kuartal menjalar menjadi inflasi, pertumbuhan ekonomi, dan pergerakan nilai tukar. Model yang digunakan mengacu pada kerangka \textit{New Keynesian} untuk ekonomi terbuka kecil (\textit{small open economy}), mengikuti Galí dan Monacelli \cite{gali2005} serta Clarida, Galí, dan Gertler \cite{clarida1999}. Model ini terdiri dari tiga persamaan inti yang saling terhubung secara rekursif. \textit{Output} hari ini dipengaruhi oleh suku bunga kemarin, inflasi hari ini dipengaruhi oleh \textit{output} hari ini, dan nilai tukar hari ini dipengaruhi oleh selisih suku bunga domestik terhadap suku bunga Amerika Serikat. Dari ketiga persamaan inti tersebut, diturunkan tiga persamaan proksi tambahan yang merepresentasikan daya beli masyarakat, neraca transaksi berjalan, dan imbal hasil obligasi pemerintah. Seluruh persamaan beroperasi pada tiap periode $t$ dari 1 sampai $T = 8$, dengan nilai awal periode $t = 0$.

\subsection{Dynamic IS Curve (Sisi Permintaan)}

Persamaan IS (\textit{Investment-Saving}) menggambarkan hubungan antara suku bunga riil dan \textit{output gap}, yaitu selisih antara \textit{output} aktual dan \textit{output} potensial suatu perekonomian.

\[
    y_t = \rho_y \, y_{t-1} - \beta \left( r_{t-1} - \pi_{t-1} - r^{*} \right) + \epsilon_t^{IS}
\]

\begin{itemize}
    \item $y_t$ adalah \textit{output gap} pada periode $t$. Jika $y_t > 0$, ekonomi sedang berjalan di atas kapasitas normalnya (\textit{overheating}). Jika $y_t < 0$, ekonomi sedang lesu di bawah potensinya.
    \item $r_{t-1} - \pi_{t-1}$ adalah suku bunga riil efektif, yaitu BI-Rate dikurangi inflasi. Inilah biaya pinjaman ``sesungguhnya'' yang dirasakan pelaku ekonomi. Semakin tinggi suku bunga riil, semakin mahal biaya kredit, sehingga investasi dan konsumsi turun, yang tercermin dari tanda negatif pada koefisien $\beta$.
    \item $r^{*}$ adalah suku bunga riil natural, yaitu tingkat suku bunga yang membuat ekonomi berada pada keseimbangan tanpa tekanan inflasi maupun resesi. Dalam spesifikasi ini $r^{*}$ disetarakan dengan target inflasi $\pi^{*} = 2,5\%$.
    \item $\rho_y$ mengukur seberapa persisten \textit{output gap}. Nilai $\rho_y = 0,60$ berarti 60\% kondisi \textit{output} periode lalu terbawa ke periode sekarang (efek inersia).
    \item $\beta$ mengukur sensitivitas \textit{output} terhadap perubahan suku bunga riil. Nilai $\beta = 0,30$ berarti kenaikan suku bunga riil sebesar satu poin persentase akan menurunkan \textit{output gap} sebesar 0,30 poin persentase.
    \item $\epsilon_t^{IS}$ adalah guncangan acak pada sisi permintaan, misalnya perubahan mendadak pada konsumsi rumah tangga atau investasi swasta yang tidak dapat diprediksi oleh model.
\end{itemize}

\subsection{Hybrid New Keynesian Phillips Curve (Sisi Penawaran dan Inflasi)}

Persamaan Kurva Phillips menghubungkan inflasi dengan \textit{output gap} dan depresiasi nilai tukar. Versi \textit{hybrid} yang digunakan di sini mencakup komponen \textit{backward-looking} karena pelaku ekonomi di Indonesia cukup signifikan dalam membentuk ekspektasi inflasi berdasarkan inflasi periode lalu.

\[
    \pi_t = \rho_\pi \, \pi_{t-1} + \kappa \, y_t + \phi \, \Delta e_t + \epsilon_t^{PC}
\]

\begin{itemize}
    \item $\pi_t$ adalah inflasi pada periode $t$, dinyatakan dalam poin persentase.
    \item $\rho_\pi$ menangkap persistensi inflasi. Nilai $\rho_\pi = 0,85$ berarti inflasi cenderung ``lengket'' dan tidak berubah drastis antar periode, sebuah karakteristik empiris perekonomian Indonesia yang memiliki komponen ekspektasi adaptif cukup kuat.
    \item $\kappa$ adalah kemiringan (\textit{slope}) Kurva Phillips, yang mengukur seberapa besar \textit{output gap} mendorong inflasi. Ketika ekonomi \textit{overheating} ($y_t > 0$), permintaan melebihi kapasitas produksi sehingga harga-harga naik. Nilai $\kappa = 0,11$ diestimasi dari data historis Indonesia oleh Juhro \cite{juhro2004}.
    \item $\Delta e_t$ adalah persen depresiasi nilai tukar rupiah terhadap USD. Tanda positif berarti rupiah melemah.
    \item $\phi$ adalah koefisien \textit{exchange rate pass-through}, yang mengukur seberapa besar depresiasi rupiah diteruskan ke harga barang dan jasa domestik. Nilai $\phi = 0,30$ berarti depresiasi 10\% akan menaikkan inflasi sekitar 3,0 poin persentase melalui kenaikan harga barang impor.
    \item $\epsilon_t^{PC}$ adalah guncangan acak pada sisi penawaran, misalnya lonjakan harga pangan akibat gagal panen atau kenaikan harga energi global.
\end{itemize}

\subsection{Uncovered Interest Rate Parity (Nilai Tukar)}

Persamaan UIP menjelaskan pergerakan nilai tukar berdasarkan selisih suku bunga domestik dan suku bunga luar negeri. Logika dasarnya sederhana. Jika BI-Rate lebih tinggi dari Fed Funds Rate, investor global akan memindahkan dananya ke Indonesia untuk memperoleh imbal hasil lebih tinggi, sehingga permintaan rupiah naik dan rupiah menguat.

\[
    \Delta e_t = -\left( r_t - r_t^{US} \right) + \rho_e + \epsilon_t^{FX}
\]

\begin{itemize}
    \item $r_t^{US}$ adalah suku bunga acuan Amerika Serikat (Fed Funds Rate) pada kuartal ke-$t$. Tidak seperti parameter lain yang bernilai konstan, $r_t^{US}$ mengikuti lintasan skenario yang merepresentasikan siklus pelonggaran yang berlanjut kemudian berbalik mengetat (\textit{easing-to-hiking cycle}) sepanjang dua tahun ke depan. Lintasan yang ditetapkan adalah $r_t^{US} = (3{,}375\%,\ 3{,}125\%,\ 3{,}125\%,\ 3{,}125\%,\ 3{,}375\%,\ 3{,}625\%,\ 3{,}875\%,\ 4{,}125\%)$, yang menggambarkan kelanjutan pemotongan suku bunga pada empat kuartal pertama hingga menyentuh dasar $3{,}125\%$, ditahan selama dua kuartal, lalu dinaikkan kembali pada dua kuartal terakhir seiring pemulihan ekonomi.
    \item Tanda negatif pada selisih $r_t - r_t^{US}$ mencerminkan hubungan terbalik. Semakin besar selisih positif (BI-Rate lebih tinggi), semakin besar aliran modal masuk sehingga rupiah menguat ($\Delta e_t$ negatif, atau apresiasi).
    \item $\rho_e$ adalah premi risiko dasar nilai tukar, ditetapkan netral sebesar 0.
    \item $\epsilon_t^{FX}$ adalah guncangan acak pada pasar valuta asing, misalnya gejolak sentimen global atau perubahan persepsi risiko terhadap negara berkembang.
    \item $\Delta e_t$ dinyatakan dalam persen depresiasi. Nilai tukar level $e_t$ (dalam Rp/USD) diperbarui dengan rumus $e_t = e_{t-1} \cdot (1 + \Delta e_t / 100)$.
\end{itemize}

\subsection{Tiga Persamaan Proksi}

Dari ketiga persamaan inti di atas, diturunkan tiga variabel proksi yang dibutuhkan untuk mengevaluasi \textit{constraint} dan fungsi objektif.

\textbf{Daya Beli Masyarakat.} Proksi daya beli didefinisikan sebagai selisih antara pertumbuhan ekonomi dan inflasi. Jika \textit{output} tumbuh lebih cepat dari inflasi, daya beli riil masyarakat meningkat. Konsep ini sederhana namun cukup untuk menangkap efek distributif kebijakan moneter untuk \textit{scope} tugas besar ini.

\[
    PP_t = y_t - \pi_t
\]

\textbf{Neraca Transaksi Berjalan.} Neraca transaksi berjalan (\textit{current account}) adalah catatan arus barang, jasa, dan transfer antara Indonesia dengan negara lain. Persamaan ini menangkap dua kekuatan yang berlawanan. Depresiasi rupiah ($\Delta e_t > 0$) membuat ekspor lebih murah bagi pembeli asing sehingga memperbaiki neraca transaksi berjalan, tercermin dari tanda positif pada $\alpha_1$. Sementara itu, \textit{output} domestik yang tinggi ($y_t > 0$) meningkatkan impor karena permintaan dalam negeri membesar, tercermin dari tanda negatif pada $\alpha_2$.

\[
    CA_t = \alpha_1 \, \Delta e_t - \alpha_2 \, y_t + \epsilon_t^{CA}
\]

Sebagai catatan, nilai $CA_t$ yang dihasilkan sudah dalam satuan persen terhadap PDB.

\textbf{Imbal Hasil Surat Berharga Negara.} Obligasi pemerintah Indonesia (SBN) menawarkan imbal hasil yang umumnya berada di atas BI-Rate karena investor menuntut kompensasi atas risiko fiskal. \textit{Spread} antara imbal hasil SBN dan BI-Rate tidak boleh terlalu lebar karena menandakan kekhawatiran pasar terhadap kemampuan fiskal pemerintah.

\[
    y_t^{SBN} = r_t + \rho_{fiskal} + \epsilon_t^{SBN}
\]

dengan $\rho_{fiskal}$ adalah premi risiko fiskal dasar sebesar 1,50 poin persentase.

\subsection{Guncangan Stokastik $\epsilon_t^{(\cdot)}$}

Seluruh notasi $\epsilon_t^{(\cdot)}$ pada persamaan di atas adalah variabel acak yang disimulasikan dengan distribusi normal. Guncangan ini merepresentasikan kejadian tak terduga seperti lonjakan harga minyak dunia, bencana alam yang mengganggu rantai pasok, atau perubahan sentimen investor global. Kehadiran guncangan stokastik memiliki implikasi penting bagi proses \textit{local search}. Evaluasi $J(s)$ untuk state yang sama bisa menghasilkan nilai sedikit berbeda karena realisasi guncangan yang berbeda. Untuk menjamin reproduktibilitas, seluruh peserta anda menggunakan \textit{seed} tetap (\texttt{seed = 42}), sehingga guncangan yang dialami oleh tiap kandidat state identik antar satu sama lain.

